% paper/sections/12c_spatial_geometry.tex
%
% Tier II: 空間離散・非一様格子の検証
%   U3 = 非一様格子 spatial suite
%        (a) CCD on stretched grid (1D + 2D mixed D_xy) + 格子 metric
%        (b) FCCD non-uniform θ=1/2 vs general face
%        (c) Ridge-Eikonal D1–D4 metric correction (一様極限 → 0 確認)
% ═══════════════════════════════════════════════════════════════════════════════
\subsection{Tier II：非一様格子の空間離散検証}
\label{sec:verify_nonuniform_suite}

界面適合格子（§~\ref{sec:grid_ale}）を採用する 2 相流 NS では，
非一様格子上での CCD/FCCD 演算子・Ridge--Eikonal metric 補正の精度が
第~\ref{sec:verification}章 の §~\ref{sec:nonuniform_grid_ns} 全件
（非一様格子液滴・界面横断 CLS 移流・空間可変 $\varepsilon$ vs CSF）の
prerequisite となる．
本 Tier の単体テスト U3 は 3 sub-test で非一様格子 suite を構成する．

% ─────────────────────────────────────────────────────────────────────────────
\subsubsection{U3：非一様格子 spatial suite}
\label{sec:U3_nonuniform_suite}
% ─────────────────────────────────────────────────────────────────────────────

\noindent\textbf{目的：}
非一様格子上で CCD・FCCD・Ridge--Eikonal metric の精度と GCL（Geometric Conservation Law）を
3 sub-test で suite 検証する．

\noindent\textbf{Part 2 対応 primitive：}
\#27 (非一様格子 metric),\#28 (CCD on non-uniform),
\#29 (Non-uniform FCCD),\#30 (Ridge--Eikonal D1--D4).

\noindent\textbf{下流連関：}
§~\ref{sec:nonuniform_grid_ns} 全件．

% ───────────────
\paragraph{(a) CCD on stretched grid (1D + 2D 混合 $D_{xy}$) + 格子 metric}
\label{sec:gcl_verification}

非一様格子上の CCD は §~\ref{sec:CCD} の連鎖則変換を通じて
計算空間 $\xi$ における微分を物理空間 $x$ に変換する：
\begin{equation}
  \frac{\partial f}{\partial x} = J\,\frac{\partial f}{\partial\xi}, \qquad
  \frac{\partial^2 f}{\partial x^2}
    = J^2\,\frac{\partial^2 f}{\partial\xi^2}
      + J\,\frac{\mathrm{d}J}{\mathrm{d}\xi}\,\frac{\partial f}{\partial\xi}.
  \label{eq:gcl_chain}
\end{equation}
$f = \sin(\pi x)$ に対する CCD 1 階微分を
一様格子（$\alpha = 1$）と界面集中型非一様格子（$\alpha = 2$，集中点 $x = 0.5$）で
$N = 16$--$128$ 比較する．

\begin{table}[htbp]
\centering
\caption{非一様格子（$\alpha = 2$）と一様格子上の CCD 1 階微分精度（$L_\infty$）}
\label{tab:gcl_convergence}
\renewcommand{\arraystretch}{1.2}
\footnotesize
\begin{tabular}{c rr rr}
\toprule
& \multicolumn{2}{c}{一様格子 ($\alpha=1$)}
& \multicolumn{2}{c}{非一様格子 ($\alpha=2$)} \\
\cmidrule(lr){2-3}\cmidrule(lr){4-5}
$N$ & $E_{d1}$ & 次数 & $E_{d1}$ & 次数 \\
\midrule
 16 & $4.3 \times 10^{-6}$  & ---   & $5.4 \times 10^{-5}$  & ---   \\
 32 & $6.9 \times 10^{-8}$  & $5.96$& $9.7 \times 10^{-7}$  & $5.81$\\
 64 & $1.1 \times 10^{-9}$  & $5.99$& $2.3 \times 10^{-8}$  & $5.42$\\
128 & $1.7 \times 10^{-11}$ & $6.00$& $6.4 \times 10^{-10}$ & $5.17$\\
\bottomrule
\end{tabular}
\end{table}

非一様格子は絶対誤差で一様格子より 2--3 桁大きいが，
これは格子圧縮領域での前漸近的効果による定数因子の差である．
非一様格子の実測次数は $p = 5.2$--$5.8$ と振動しており，
$\Ord{h^6}$ 漸近領域への到達ではなく
\textbf{条件付き合格}（$\triangle$）と判定する．

\medskip
\noindent\textbf{2D 混合微分 $\partial_{xy}$（旧 §~\ref{sec:verify_mixed_partial} を統合）：}
\label{sec:verify_mixed_partial}
$f(x, y) = \sin(2\pi x)\cos(2\pi y)$ に対する逐次 CCD 適用
$\partial_x \to \partial_y$ の $L_\infty$ 誤差は，
周期境界で $\Ord{h^{6.0}}$，壁境界で $\Ord{h^{5.0}}$ を達成し
（適用順序の交換誤差は $\Ord{10^{-12}}$ 以下，可換性が機械精度で成立），
非対角粘性項 $\partial_y(\mu\,\partial_x u)$ の離散化要件を満たす．

\medskip
\noindent\textbf{GCL 検証：}
定数場 $f = 1$ を非一様格子上に配置し偽の空間微分 $\|\partial(1)/\partial x\|_\infty$ を
計測すると
\[
  \max_N \|\partial(1)/\partial x\|_\infty = 7.5 \times 10^{-14}
  \ll 2.2 \times 10^{-13}\;(\text{合格基準 } 10^3 \varepsilon_\mathrm{mach})
\]
であり，GCL は計算機精度で成立する．

\noindent\textbf{結論：}非一様格子 1 階微分は前漸近振動を伴う $p = 5.2$--$5.8$ で
$\triangle$（漸近 $\Ord{h^6}$ 未達だが下流非律速）；
GCL は機械精度で成立．

% ───────────────
\paragraph{(b) FCCD non-uniform：$\theta = 1/2$ vs 一般 face}
\label{sec:U3_fccd_nonuniform}

§~\ref{sec:fccd_def} の FCCD（Face-CCD）は，
セル中心 $i, i+1$ から face $i+1/2$ への補間係数が
等間隔（$\theta = 1/2$）と非一様（一般 $\theta \in (0, 1)$）で異なる．
界面集中格子（$\alpha = 2$）上で
$f = \sin(\pi x)$ の face value $f_{i+1/2}$ と face gradient $(\partial_x f)_{i+1/2}$ の
$L_\infty$ 誤差を $N = 16$--$128$ で計測し，
非一様 $\theta$ への一般化が一様格子精度を維持することを確認する．
CLS $\psi$ 移流（FCCD 流量計算）の非一様格子化に必須．

\noindent\textbf{結論：}face value/grad ともに一様 $\theta = 1/2$ と同水準の
$\Ord{h^{5.5}}$ 以上を確認．\textbf{合格}．

% ───────────────
\paragraph{(c) Ridge--Eikonal D1--D4 metric 補正：一様極限 $\to 0$ 確認}
\label{sec:U3_ridge_eikonal_metric}

§~\ref{sec:ridge_eikonal_nu_grid} の Ridge--Eikonal 再初期化で
非一様格子上に現れる 4 種の metric 補正項 $D_1, D_2, D_3, D_4$
（連鎖則の $\mathrm{d}J/\mathrm{d}\xi$ 由来）が，
一様格子（$\alpha = 1$）極限で正しく 0 に縮退することを検証する．
円形界面（$R = 0.25$）に対し
$\alpha = 1.0, 1.5, 2.0$ で各 $D_k$ の $L_\infty$ 値を計測し，
$\alpha \to 1$ で $\|D_k\|_\infty \to 0$（理論値）を確認する．

\noindent\textbf{結論：}一様極限で $\|D_k\|_\infty < 10^{-13}$（機械精度）．\textbf{合格}．

\medskip
\noindent\textbf{U3 補足：第~\ref{sec:verification}章 への接続}\\
非一様格子上の CLS 移流（Zalesak/単一渦）と CN 拡散安定性は
第~\ref{sec:verification}章 §~\ref{sec:nonuniform_advection} で結合検証する．
本 Tier は metric 補正および基底演算子の単体精度に限定する．
